
\subsubsection{6.1 Key Findings}

\textbf{Finding 1: The graph coordinate system matters more than the symmetry group for cross-architecture SSL.} The ~220\% improvement over SANE arises primarily from graph-based encoding, which provides architecture-agnostic relational structure. Given graph-based encoding, permutation-only equivariance outperforms scale+permutation equivariance for ViT transfer ($\Delta R^{2}=+0.021$ under comparable training conditions in h-m1; $\Delta R^{2}=+0.143$ in the h-m2 ablation, though this value is partially confounded). This suggests the appropriate inductive bias is architecture-normalization-dependent: permutation equivariance for LayerNorm architectures, potentially scale+permutation for BatchNorm or post-ReLU architectures without adaptive normalization.

\textbf{Finding 2: MMD is an unreliable metric for cross-architecture transfer evaluation.} SANE achieves lower MMD than graph-based encoders despite substantially worse downstream $R^2$. Representation collapse produces trivially low MMD; this is a measurement artifact, not evidence of alignment. For cross-architecture property prediction, $R^2$ from a linear probe is a more reliable evaluation metric. We additionally note that different experiments in this work measured different quantities under the label ''MMD'' (cross-architecture MMD in h-e1 vs. within-zoo subpopulation MMD in h-m2); these quantities are not directly comparable.

\textbf{Finding 3: Property prediction and functional model generation are separable.} The latent interpolation failure (acc=0.100) was caused by a decoder design limitation --- the decoder reconstructed edge statistics rather than functional weights --- not by a failure of the latent representation to encode model properties (which $R^{2}=0.231$ demonstrates). A dedicated hypernetwork-style decoder with architecture-aware output projections would be required for functional interpolation.

\subsubsection{6.2 Limitations}

\textbf{Statistical power.} All cross-architecture $R^2$ results are from a single training seed with 53 ViT-S/16 models. With n=1 seed, statistical significance is not assessable --- t-test p-values computed in h-m1 are degenerate (p=0.5 for n=1). Magnitude estimates carry high uncertainty. Full evaluation with 5 seeds and 250+ ViT models is required for reproducible statistical claims.

\textbf{Ablation fairness.} Table 2 (h-m2) compares EquiSSL from a 50-epoch checkpoint against EquiSSL-perm from a 100-epoch checkpoint. The observed $\Delta R^{2}=+0.143$ is a confounded estimate. The h-m1 comparison provides a more controlled estimate ($\Delta R^{2}=+0.021)$, but even this uses a shared probe setup rather than a strict train-from-scratch ablation with identical epoch counts and initialization seeds.

\textbf{Causal status of the LayerNorm gauge hypothesis.} The hypothesis is empirically motivated and consistent with observations but not directly tested as a causal mechanism. Confirmation requires evaluating the same scale vs. permutation-only comparison on non-LayerNorm target architectures (e.g., ReLU MLP, BatchNorm CNN), which was not performed in this work.

\textbf{Training data scale.} Approximately 3,000 of roughly 30,000 available SANE MultiZoo CNN checkpoints were used for encoder training. The effect of training data scale on cross-architecture transfer is unknown.

\textbf{ViT test set size.} 53 of 250 available ViT-S/16 models were used for evaluation, due to local availability constraints. The representativeness of this subset relative to the full zoo is not characterized.

\textbf{Decoder design.} The graph decoder in h-m3 cannot generate functional weight tensors. A hypernetwork-style decoder with architecture-aware output projections is the natural correction.

\textbf{MMD comparability.} Cross-architecture MMD was computed only for SANE and EquiSSL in h-e1. No cross-architecture MMD is available for EquiSSL-perm, limiting the scope of the RQ2 analysis.

\textbf{Unverified citation.} [Anonymous2025LayerNorm; arXiv:2510.08300] could not be confirmed in Semantic Scholar at the time of submission. Claims referencing this citation are hedged accordingly throughout the paper.

\subsubsection{6.3 Broader Impact}

Weight-space SSL methods that generalize across architecture families could reduce the cost of model zoo analysis --- enabling property prediction for arbitrary pre-trained models without dedicated training data for each new architecture. This could facilitate model evaluation and trustworthy deployment at scale. Weight-space representations could theoretically reveal information about training data from model checkpoints; more capable future methods may raise more significant concerns along this dimension. Code and model checkpoints are released to enable reproducibility and independent evaluation.

